% GENERATED FILE. Edit canonical lessons, then run npm run generate:books.
% canonical-source-hash: fnv1a64:a52fbcc16fbe999a
% canonical-lessons: TE-C399-kartavyam, TE-C399-sveccha, TE-C399-samanatvam, TE-C399-anyayam, TE-C399-gauravincu

\chapter{Rights and Duties}
\label{ch:te-rights-and-duties}
\hlchaptermodality{399}

\begin{chapteropening}
\textbf{By the end of this chapter:} \emph{I can talk about duty, freedom, equality and justice.}

Complete the last lesson of chapter 399: I can talk about duty, freedom, equality and justice.
\end{chapteropening}

\section[\texorpdfstring{\te{కర్తవ్యం} (kartavyaṁ)}{కర్తవ్యం (kartavyaṁ)}]{\te{కర్తవ్యం} (kartavyaṁ) — a duty; an obligation}

\label{lesson:TE-C399-kartavyam}



\begin{quote}
Before the new one: say the Telugu for a punishment, then the Telugu for to ban.
\end{quote}

\subsection*{You'll want to know: \te{కర్తవ్యం}}
\textbf{\te{కర్తవ్యం}} — \emph{kartavyaṁ} — \textquotedblleft{}a duty; an obligation\textquotedblright{}.

\begin{quote}
\te{తల్లిదండ్రులను} \te{చూసుకోవడం} \te{మన} \te{కర్తవ్యం}.
\end{quote}

\begin{quote}
\emph{tallidaṇḍrulanu cūsukōvaḍaṁ mana kartavyaṁ.}
\end{quote}

\begin{quote}
Looking after our parents is our duty.
\end{quote}

A higher word than \te{బాధ్యత}: a moral duty.

\begin{cousinweb}
\textbf{\te{కర్తవ్యం}} (\emph{kartavyaṁ}) is borrowed from Sanskrit \emph{kartavya}, \textquotedblleft{}to be done\textquotedblright{}, and as a noun \textquotedblleft{}that which ought to be done, a duty\textquotedblright{} (Monier-Williams), from \emph{kṛ}, \textquotedblleft{}to do\textquotedblright{}.
\end{cousinweb}

\begin{grammarlens}[title={What you've built}]
One more word for giving reasons and opinions.
\end{grammarlens}

\subsection*{Your turn}
\begin{itemize}
\raggedright
  \item \emph{Say it:} \emph{kartavyaṁ}
  \item \emph{Say it:} \emph{kartavyaṁ}, once more
  \item \emph{Say it:} \emph{tallidaṇḍrulanu cūsukōvaḍaṁ mana kartavyaṁ.} — Looking after our parents is our duty.
  \item \emph{Recall:} say the Telugu for a punishment, then the Telugu for to ban, then say \emph{kartavyaṁ} again
\end{itemize}

\begin{tcolorbox}[breakable,colback=teal!4,colframe=teal!35!black,title={Before you move on}]
What does \te{కర్తవ్యం} mean? (\textquotedblleft{}A duty; an obligation\textquotedblright{}.) Say it once more.
\end{tcolorbox}
\section[\texorpdfstring{\te{స్వేచ్ఛ} (svēccha)}{స్వేచ్ఛ (svēccha)}]{\te{స్వేచ్ఛ} (svēccha) — freedom; liberty}

\label{lesson:TE-C399-sveccha}



\begin{quote}
Before the new one: say the Telugu for to ban, then the Telugu for a duty.
\end{quote}

\subsection*{You'll want to know: \te{స్వేచ్ఛ}}
\textbf{\te{స్వేచ్ఛ}} — \emph{svēccha} — \textquotedblleft{}freedom; liberty\textquotedblright{}.

\begin{quote}
\te{ప్రతి} \te{మనిషికి} \te{మాట్లాడే} \te{స్వేచ్ఛ} \te{ఉండాలి}.
\end{quote}

\begin{quote}
\emph{prati maniṣiki māṭlāḍē svēccha uṇḍāli.}
\end{quote}

\begin{quote}
Every person should have the freedom to speak.
\end{quote}

\te{స్వ}, one’s own, + \te{ఇచ్ఛ}, wish. \te{స్వేచ్ఛగా}, freely.

\begin{cousinweb}
\textbf{\te{స్వేచ్ఛ}} (\emph{svēccha}) is borrowed from Sanskrit \emph{svecchā}, \textquotedblleft{}one’s own wish, free will\textquotedblright{} (Monier-Williams), from \emph{sva}, \textquotedblleft{}own\textquotedblright{}, and \emph{icchā}, \textquotedblleft{}wish\textquotedblright{}.
\end{cousinweb}

\begin{grammarlens}[title={What you've built}]
One more word for giving reasons and opinions.
\end{grammarlens}

\subsection*{Your turn}
\begin{itemize}
\raggedright
  \item \emph{Say it:} \emph{svēccha}
  \item \emph{Say it:} \emph{svēccha}, once more
  \item \emph{Say it:} \emph{prati maniṣiki māṭlāḍē svēccha uṇḍāli.} — Every person should have the freedom to speak.
  \item \emph{Recall:} say the Telugu for to ban, then the Telugu for a duty, then say \emph{svēccha} again
\end{itemize}

\begin{tcolorbox}[breakable,colback=teal!4,colframe=teal!35!black,title={Before you move on}]
What does \te{స్వేచ్ఛ} mean? (\textquotedblleft{}Freedom; liberty\textquotedblright{}.) Say it once more.
\end{tcolorbox}
\section[\texorpdfstring{\te{సమానత్వం} (samānatvaṁ)}{సమానత్వం (samānatvaṁ)}]{\te{సమానత్వం} (samānatvaṁ) — equality}

\label{lesson:TE-C399-samanatvam}



\begin{quote}
Before the new one: say the Telugu for a duty, then the Telugu for freedom.
\end{quote}

\subsection*{You'll want to know: \te{సమానత్వం}}
\textbf{\te{సమానత్వం}} — \emph{samānatvaṁ} — \textquotedblleft{}equality\textquotedblright{}.

\begin{quote}
\te{ఆడవాళ్ళకు} \te{మగవాళ్ళకు} \te{సమానత్వం} \te{ఉండాలి}.
\end{quote}

\begin{quote}
\emph{āḍavāḷḷaku magavāḷḷaku samānatvaṁ uṇḍāli.}
\end{quote}

\begin{quote}
There should be equality between women and men.
\end{quote}

\te{సమాన}, equal, + -\te{త్వం}, -ness. \te{సమానంగా}, equally.

\begin{cousinweb}
\textbf{\te{సమానత్వం}} (\emph{samānatvaṁ}) is borrowed from Sanskrit \emph{samānatva}, \textquotedblleft{}equality\textquotedblright{} (Monier-Williams), from \emph{samāna}, \textquotedblleft{}same, alike\textquotedblright{} (CDIAL 13211), and the abstract ending \emph{-tva}.
\end{cousinweb}

\begin{grammarlens}[title={What you've built}]
One more word for giving reasons and opinions.
\end{grammarlens}

\subsection*{Your turn}
\begin{itemize}
\raggedright
  \item \emph{Say it:} \emph{samānatvaṁ}
  \item \emph{Say it:} \emph{samānatvaṁ}, once more
  \item \emph{Say it:} \emph{āḍavāḷḷaku magavāḷḷaku samānatvaṁ uṇḍāli.} — There should be equality between women and men.
  \item \emph{Recall:} say the Telugu for a duty, then the Telugu for freedom, then say \emph{samānatvaṁ} again
\end{itemize}

\begin{tcolorbox}[breakable,colback=teal!4,colframe=teal!35!black,title={Before you move on}]
What does \te{సమానత్వం} mean? (\textquotedblleft{}Equality\textquotedblright{}.) Say it once more.
\end{tcolorbox}
\section[\texorpdfstring{\te{అన్యాయం} (anyāyaṁ)}{అన్యాయం (anyāyaṁ)}]{\te{అన్యాయం} (anyāyaṁ) — injustice; a wrong}

\label{lesson:TE-C399-anyayam}



\begin{quote}
Before the new one: say the Telugu for freedom, then the Telugu for equality.
\end{quote}

\subsection*{You'll want to know: \te{అన్యాయం}}
\textbf{\te{అన్యాయం}} — \emph{anyāyaṁ} — \textquotedblleft{}injustice; a wrong\textquotedblright{}.

\begin{quote}
\te{పేదవాళ్ళకు} \te{అన్యాయం} \te{జరిగింది}.
\end{quote}

\begin{quote}
\emph{pēdavāḷḷaku anyāyaṁ jarigindi.}
\end{quote}

\begin{quote}
An injustice was done to the poor.
\end{quote}

\te{అ}-, not, + \te{న్యాయం}, justice; \te{న్యాయమైన}, fair, which you know, is the same word.

\begin{cousinweb}
\textbf{\te{అన్యాయం}} (\emph{anyāyaṁ}) is borrowed from Sanskrit \emph{anyāya}, \textquotedblleft{}unjust or unlawful action\textquotedblright{} (Monier-Williams), the opposite of \emph{nyāya}, \textquotedblleft{}rule, justice\textquotedblright{}.
\end{cousinweb}

\begin{grammarlens}[title={What you've built}]
One more word for giving reasons and opinions.
\end{grammarlens}

\subsection*{Your turn}
\begin{itemize}
\raggedright
  \item \emph{Say it:} \emph{anyāyaṁ}
  \item \emph{Say it:} \emph{anyāyaṁ}, once more
  \item \emph{Say it:} \emph{pēdavāḷḷaku anyāyaṁ jarigindi.} — An injustice was done to the poor.
  \item \emph{Recall:} say the Telugu for freedom, then the Telugu for equality, then say \emph{anyāyaṁ} again
\end{itemize}

\begin{tcolorbox}[breakable,colback=teal!4,colframe=teal!35!black,title={Before you move on}]
What does \te{అన్యాయం} mean? (\textquotedblleft{}Injustice; a wrong\textquotedblright{}.) Say it once more.
\end{tcolorbox}
\section[\texorpdfstring{\te{గౌరవించు} (gauraviñcu)}{గౌరవించు (gauraviñcu)}]{\te{గౌరవించు} (gauraviñcu) — to honour; to respect}

\label{lesson:TE-C399-gauravincu}



\begin{quote}
Before the new one: say the Telugu for equality, then the Telugu for injustice.
\end{quote}

\subsection*{You'll want to know: \te{గౌరవించు}}
\textbf{\te{గౌరవించు}} — \emph{gauraviñcu} — \textquotedblleft{}to honour; to respect\textquotedblright{}.

\begin{quote}
\te{మనం} \te{పెద్దలను} \te{గౌరవించాలి}.
\end{quote}

\begin{quote}
\emph{manaṁ peddalanu gauraviñcāli.}
\end{quote}

\begin{quote}
We should respect our elders.
\end{quote}

Past (I …): \te{గౌరవించాను}, \emph{gauraviñcānu}.

The verb of \te{గౌరవం}, respect, which you know.

\begin{cousinweb}
\textbf{\te{గౌరవించు}} (\emph{gauraviñcu}) is built on Sanskrit \emph{gaurava}, \textquotedblleft{}weight\textquotedblright{}, and so \textquotedblleft{}venerableness\textquotedblright{} (CDIAL 4346), with the Telugu verb ending -\te{ఇంచు}.
\end{cousinweb}

\begin{grammarlens}[title={What you've built}]
One more word for giving reasons and opinions.
\end{grammarlens}

\subsection*{Your turn}
\begin{itemize}
\raggedright
  \item \emph{Say it:} \emph{gauraviñcu}
  \item \emph{Say it:} \emph{gauraviñcu}, once more
  \item \emph{Say it:} \emph{manaṁ peddalanu gauraviñcāli.} — We should respect our elders.
  \item \emph{Recall:} say the Telugu for equality, then the Telugu for injustice, then say \emph{gauraviñcu} again
\end{itemize}

\begin{tcolorbox}[breakable,colback=teal!4,colframe=teal!35!black,title={Before you move on}]
What does \te{గౌరవించు} mean? (\textquotedblleft{}To honour; to respect\textquotedblright{}.) Say it once more.
\end{tcolorbox}
